\section{Resultados}

\subsection{Validação cruzada agrupada}

A validação cruzada agrupada em cinco folds constitui o resultado primário. No
limiar 0,50 e sem TTA, a ConvNeXt-Tiny obteve AUC ROC de
$0{,}9709\pm0{,}0059$ e AUC PR de $0{,}9607\pm0{,}0121$. A
Tabela~\ref{tab:cv-media} resume as métricas agregadas e a
Figura~\ref{fig:cv-folds} mostra que a variação principal ocorre nas métricas
dependentes do limiar, não no ranqueamento global.

\begin{table}[ht]
\centering
\caption{Resultado primário da validação cruzada em cinco folds agrupada por
imagem, com limiar 0,50 e sem TTA. Valores em média $\pm$ desvio-padrão.}
\label{tab:cv-media}
\begin{tabular}{lr}
\toprule
Métrica & Resultado \\
\midrule
AUC ROC & $0{,}9709\pm0{,}0059$ \\
AUC PR & $0{,}9607\pm0{,}0121$ \\
Acurácia balanceada & $0{,}9091\pm0{,}0115$ \\
Sensibilidade & $0{,}8994\pm0{,}0480$ \\
Especificidade & $0{,}9188\pm0{,}0400$ \\
\emph{F1} & $0{,}8926\pm0{,}0128$ \\
MCC & $0{,}8189\pm0{,}0208$ \\
Kappa & $0{,}8164\pm0{,}0223$ \\
\emph{Brier score} & $0{,}0699\pm0{,}0079$ \\
\bottomrule
\end{tabular}
\end{table}

\begin{figure}[ht]
\centering
\begin{tikzpicture}
\begin{axis}[
 width=.72\textwidth,height=4.6cm,
 kariellyaxis,
 ybar,bar width=9pt,
 symbolic x coords={1,2,3,4,5},xtick=data,
 xlabel={Fold},ylabel={M\'etrica dependente do limiar},
 ymin=.80,ymax=.99,
 legend style={at={(.5,-.20)},anchor=north,legend columns=3}
]
\addplot[fill=black!70,draw=black!70] table[x=fold,y=sensitivity,col sep=comma]{dados/cv_folds_plot.csv};
\addlegendentry{Sens.}
\addplot[fill=black!20,draw=black!45] table[x=fold,y=specificity,col sep=comma]{dados/cv_folds_plot.csv};
\addlegendentry{Espec.}
\addplot[black,thick,mark=*,sharp plot] table[x=fold,y=balanced_acc,col sep=comma]{dados/cv_folds_plot.csv};
\addlegendentry{Acc. bal.}
\end{axis}
\end{tikzpicture}
\caption{Métricas dependentes do limiar nos cinco folds agrupados. O fold 1
privilegia especificidade e perde sensibilidade; os folds 2 e 3 mostram o
comportamento oposto.}
\label{fig:cv-folds}
\end{figure}

A Figura~\ref{fig:cv-folds} justifica reportar a validação cruzada além da AUC:
no fold 1, a sensibilidade cai para 0,8253 enquanto a especificidade chega a
0,9681; nos folds 2 e 3, a recuperação de anormais aumenta, mas com menor
especificidade. Portanto, a variação relevante para uso em triagem está no
limiar operacional e na composição do fold, não apenas na capacidade de
ranqueamento do modelo.

\subsection{Análise detalhada em uma partição}

Na partição retida, a inferência da ConvNeXt-Tiny com TTA produziu AUC ROC de
0,9832 e AUC PR de 0,9827. A temperatura estimada na validação foi
$T=1{,}210$ e o \emph{Brier score} foi 0,0473. O ECE passou de 0,0247 para
0,0130 após o escalonamento, indicando melhora de calibração nessa partição.
Como a temperatura é monotônica, ela não altera as AUCs; a pequena mudança de
ranqueamento em relação à inferência sem TTA decorre do TTA.

\subsection{Efeito do limiar operacional}

Os três regimes evidenciam a troca entre falsos negativos e falsos positivos.
O limiar 0,50 apresentou o maior MCC e a maior acurácia balanceada. O limiar de
Youden aumentou a sensibilidade para 0,9522 com especificidade 0,9283. O regime
de alta sensibilidade atingiu 0,9701 de sensibilidade, mas reduziu a
especificidade para 0,8723, ponto aceitável apenas se a prioridade operacional
for reduzir omissões e encaminhar mais casos para revisão humana. A
Tabela~\ref{tab:confusao-limiares} explicita essa troca em contagens absolutas:
ao reduzir o limiar, os falsos negativos caem de 60 para 30, enquanto os falsos
positivos sobem de 64 para 146.

\begin{table}[ht]
\centering
\caption{Matriz de decisão na partição retida para três regimes de limiar.}
\label{tab:confusao-limiares}
\scriptsize
\setlength{\tabcolsep}{3.5pt}
\begin{tabular}{lrrrrrrr}
\toprule
Regime & Limiar & TN & FP & FN & TP & Sens. & Espec. \\
\midrule
0,50 & 0,500 & 1079 & 64 & 60 & 944 & 0,9402 & 0,9440 \\
Youden & 0,287 & 1061 & 82 & 48 & 956 & 0,9522 & 0,9283 \\
Alta sens. & 0,099 & 997 & 146 & 30 & 974 & 0,9701 & 0,8723 \\
\bottomrule
\end{tabular}
\end{table}

\subsection{Comparação com ResNet50}

A comparação interna com ResNet50 foi realizada na mesma partição retida. A
ConvNeXt-Tiny apresentou maior capacidade de ranqueamento, com AUC ROC 0,9816 e
AUC PR 0,9805, contra 0,9710 e 0,9750 da ResNet50. A ResNet50, por outro lado,
obteve maior especificidade, acurácia balanceada e MCC no limiar avaliado. Esse
resultado indica que a vantagem da ConvNeXt-Tiny não é absoluta: ela ordena
melhor as amostras, enquanto a ResNet50 foi competitiva e mais favorável em
algumas métricas dependentes do limiar. A comparação deve ser lida como
exploratória, pois foi realizada apenas na partição retida, e não nos cinco
folds agrupados.

\begin{table}[ht]
\centering
\caption{Comparação entre ConvNeXt-Tiny e ResNet50 na partição retida.}
\label{tab:baseline}
\kariellytablefont
\kariellytablecols
\begin{tabular}{lrrrrrrr}
\toprule
Modelo & AUC ROC & AUC PR & Acc. bal. & Sens. & Espec. & F1 & MCC \\
\midrule
ConvNeXt-Tiny & \textbf{0,9816} & \textbf{0,9805} & 0,9371 & \textbf{0,9373} &
0,9370 & 0,9331 & 0,8738 \\
ResNet50 & 0,9710 & 0,9750 & \textbf{0,9401} & 0,9064 & \textbf{0,9738} &
\textbf{0,9362} & \textbf{0,8852} \\
\bottomrule
\end{tabular}
\end{table}

\subsection{Representatividade e erros por categoria Bethesda}

A composição da partição retida favorece categorias mais facilmente separáveis:
HSIL corresponde a 52,7\% das células anormais no teste, ante 35,8\% no
conjunto completo, enquanto LSIL corresponde a 16,4\%, ante 28,6\% no conjunto
completo. Essa super-representação de HSIL e sub-representação de LSIL ajudam a
explicar por que os resultados da partição única são superiores às médias da
validação cruzada.

No limiar de Youden, a Tabela~\ref{tab:bethesda} mostra 92,8\% de acerto para
NILM, 81,7\% para ASC-US, 94,5\% para LSIL, 90,2\% para ASC-H, 99,4\% para HSIL
e 100\% para SCC. O resultado de ASC-US, baseado em 115 células dessa partição,
sugere uma dificuldade relevante, mas não sustenta sozinho uma conclusão
populacional. SCC também possui suporte reduzido, com 42 células.

\begin{table}[ht]
\centering
\caption{Acerto binário por categoria Bethesda na partição retida, usando o
limiar de Youden (0,287).}
\label{tab:bethesda}
\small
\setlength{\tabcolsep}{5pt}
\begin{tabular}{lrrlr}
\toprule
Categoria & $n$ & Acerto (\%) & IC95\% Wilson & $P$(anormal) média \\
\midrule
NILM & 1143 & 92,8 & 91,2--94,2\% & 0,091 \\
ASC-US & 115 & 81,7 & 73,7--87,7\% & 0,741 \\
LSIL & 165 & 94,5 & 90,0--97,1\% & 0,898 \\
ASC-H & 153 & 90,2 & 84,5--94,0\% & 0,860 \\
HSIL & 529 & 99,4 & 98,3--99,8\% & 0,956 \\
SCC & 42 & 100,0 & 91,6--100,0\% & 0,970 \\
\bottomrule
\end{tabular}
\end{table}
